% File: part2/chapters3/intro.tex

\chapter{TỐI ƯU HÓA, TRIỂN KHAI VÀ VẬN HÀNH (MLOps)}
\label{chap:mlops}

Hành trình của một mô hình học máy không kết thúc ở epoch huấn luyện cuối cùng hay ở con số 95\% F1-score trên tập kiểm tra. Giá trị thực sự của nó chỉ được hiện thực hóa khi nó được đưa vào một sản phẩm, phục vụ hàng nghìn, thậm chí hàng triệu người dùng, và hoạt động một cách đáng tin cậy 24/7. Chào mừng bạn đến với thế giới của MLOps -- Vận hành Học máy.

Trong chương này, chúng ta sẽ vượt ra khỏi khuôn khổ của sổ tay Jupyter và đối mặt với những thách thức của thế giới thực. Làm thế nào để phục vụ một mô hình 70 tỷ tham số với độ trễ thấp? Làm thế nào để đóng gói một ứng dụng AI phức tạp để nó có thể chạy ở bất cứ đâu? Làm thế nào để giám sát “sức khỏe” của mô hình sau khi triển khai và biết khi nào cần huấn luyện lại nó?

Chúng ta sẽ khám phá các thành phần cơ sở hạ tầng quan trọng như Cơ sở dữ liệu Vector, các thư viện tối ưu hóa mô hình, các framework phục vụ chuyên dụng, và cuối cùng là phác thảo một vòng đời MLOps hoàn chỉnh cho LLM. Đây là những kỹ năng sẽ biến bạn từ một nhà khoa học dữ liệu thành một kỹ sư học máy toàn diện.